\begin{abstract}
Autonomous language agents frequently suffer from planning-quality degradation and compounding error propagation when executing long-horizon, multi-step workflows. While intermediate verification and state rollback can intercept corruptions before trajectory completion, verification calls incur significant computational, latency, and monetary costs. Under realistic resource constraints where an orchestrator can allocate only a small number of verification checks ($K \ll H$), the critical architectural question is: \emph{where should verification checkpoints be placed?}

In this paper, we present the first controlled comparative study evaluating verification scheduling strategies under strictly matched verification budgets. Across 90 Directed Acyclic Graph (DAG) task environments evaluating 4,680 execution episodes under controlled fault injection, we compare four conditions: (A) Final-Only verification ($K=1$), (B) Fixed-Periodic verification ($K=3$), (C) Dependency-Milestone verification placed at multi-parent join boundaries ($K=3$), and (D) Every-Step verification ($K=H$). Under strictly matched realized verification budgets ($K=2.7436 \pm 0.5872$ calls across both Arms B and C, $p = 1.0000$), our empirical findings decisively reject the intuitive hypothesis that dependency integration milestones provide superior error containment. Milestone verification permitted a statistically significant $3.91\%$ increase in the Downstream Taint-Consumption Fraction (TCF) compared to fixed-periodic verification ($\mu_{\text{TCF}} = 0.9005$ vs. $0.8614$; paired permutation $p < 0.0001$; 95\% bootstrap CI $[+0.0366, +0.0414]$; Cohen's $d = 3.35$) and incurred a $2.31$-step greater fault detection delay ($7.81$ vs. $5.50$ steps). We demonstrate that this underperformance stems from the \emph{structural clumping phenomenon}: because parallel processing branches must finish before joining, dependency milestones naturally cluster in the middle and late stages of long-horizon graphs, leaving early execution branches unmonitored. Furthermore, we show that intermediate verification decouples error containment from final task completion: under single-replay recovery, completion rates remained identical ($38.46\%$) across all arms, but periodic verification reduced recovery re-execution waste by $38.9\%$ relative to final-only checking. These results provide actionable architectural principles for designing cost-effective verification harnesses in compound AI systems.
\end{abstract}
